\section{Premise}\label{sec:premise}

\subsection{Motivation: bubbles in rocket turbopumps}

A liquid rocket engine needs its propellants delivered at high pressure, and the turbopump
that does this is one of the hardest parts of the engine to design. To keep the propellant
tanks light, the tank pressure is kept low, which means the liquid arrives at the pump with
only a small margin above its vapor pressure. The first stage of the pump, the
\emph{inducer}, is an axial impeller built to run with some vapor present. Near its blade
leading edges and tip clearance the local pressure drops below the vapor pressure, and small
gas and vapor nuclei in the liquid grow into cavitation bubbles
\cite{brennen1994pumps,franc2004}.

Those bubbles are then carried downstream into higher pressure, and they collapse. A
collapsing bubble focuses the energy of the surrounding liquid into a tiny volume over a
very short time, producing high wall speeds and large local pressures. Repeated near a
surface, this pits and erodes blades. At the pump level, enough cavitation degrades the head
the pump delivers and can drive unsteady behavior such as rotating cavitation and cavitation
surge \cite{brennen1994pumps}. The same physics shows up in propellant feed lines, valves and
injectors.

Real inducer cavitation involves clouds and sheets of bubbles, walls and non-spherical
collapse. We study the simplest building block: a single spherical bubble in a large body of
liquid whose far-field pressure changes in time. The Rayleigh--Plesset equation describes
this bubble \cite{rayleigh1917,plesset1949,brennen1995}. It is simple enough to study
carefully and still captures the violent, strongly nonlinear collapse that makes cavitation
damaging. We use water at 20\,$^\circ$C as the working liquid because its properties are
well known and the literature for checking our results is largest. Cryogenic propellants
behave differently, partly because evaporation and condensation cool and heat the bubble
(the ``thermal effect'' that suppresses cavitation in cryogens \cite{brennen1995}), and our
model leaves that out.

\subsection{Research questions}\label{sec:questions}

\begin{enumerate}
  \item[\textbf{Q1.}] \textbf{How hard does a bubble collapse after a pressure jump?}
  A nucleus that has settled to equilibrium in the low-pressure region near a blade is
  suddenly exposed to a higher recovery pressure. How do the minimum radius
  $R_{\min}/R_0$, the peak inward wall speed $|\dot R|_{\max}$ and the peak gas pressure
  $p_{g,\max}$ scale with the pressure ratio $\Pi = p_1/p_0$ and the initial radius $R_0$?
  At what pressure ratio does the gas inside stop mattering, so that the collapse looks like
  the empty-cavity Rayleigh limit until the very end?

  \item[\textbf{Q2.}] \textbf{How does a bubble respond to oscillating pressure, and where does it
  stop behaving linearly?} With the ambient pressure oscillating at frequency $\omega$ and
  relative amplitude $\varepsilon$, where is the resonance peak, and how far does it shift
  from the small-amplitude natural frequency $\omega_0$ as $\varepsilon$ grows? At what
  amplitude does the response turn from gentle oscillation into large growth followed by
  violent collapse?

  \item[\textbf{Q3.}] \textbf{Which physics controls the collapse, and when?} Gas compression,
  viscosity and surface tension all resist collapse. Over the range of $R_0$ and $\Pi$ we
  sweep, which one sets $R_{\min}$ and the peak pressure? We answer this by switching terms
  off one at a time and comparing against the dimensionless groups of
  Section~\ref{sec:nondim}, producing a regime map.
\end{enumerate}

Q1 and Q2 are the core of the project; Q3 is a stretch goal that reuses the same runs.
None of these questions can be answered from a closed-form solution, as explained in
Section~\ref{sec:whynumerics}.

\subsection{The initial value problem}\label{sec:ivp}

The bubble radius $R(t)$ obeys the Rayleigh--Plesset equation
\cite{plesset1949,brennen1995,plesset1977}
\begin{equation}\label{eq:rp}
  R\ddot R + \frac{3}{2}\dot R^2
  = \frac{1}{\rho}\left[ p_v + \pgz\left(\frac{R_0}{R}\right)^{3\kappa}
    - \pinf(t) - \frac{2\sigma}{R} - \frac{4\mu\dot R}{R} \right],
  \qquad R(0) = R_0,\quad \dot R(0) = 0,
\end{equation}
where dots are time derivatives. The bracket is the pressure in the liquid just outside the
bubble wall minus the far-field pressure. The initial gas pressure is set so the bubble
starts in equilibrium with the ambient pressure at $t = 0$,
\begin{equation}\label{eq:pg0}
  \pgz = \pinf(0) - p_v + \frac{2\sigma}{R_0},
\end{equation}
except in Case A, where the bubble is empty ($\pgz = 0$). With the state
$\mathbf{y} = [R,\ \dot R]^{\mathsf T}$, Eq.~\eqref{eq:rp} becomes the first-order system
\begin{equation}\label{eq:system}
  \frac{d}{dt}\begin{bmatrix} R \\ \dot R \end{bmatrix}
  = \begin{bmatrix}
      \dot R \\[2pt]
      \dfrac{1}{R}\left\{ \dfrac{1}{\rho}\left[ p_v + \pgz\left(\dfrac{R_0}{R}\right)^{3\kappa}
      - \pinf(t) - \dfrac{2\sigma}{R} - \dfrac{4\mu\dot R}{R} \right] - \dfrac{3}{2}\dot R^2 \right\}
    \end{bmatrix},
\end{equation}
which is what our integrators advance.

\paragraph{Symbols.} $R$ is the bubble radius and $R_0$ its initial value; $\rho$, $\mu$ and
$\sigma$ are the liquid density, dynamic viscosity and surface tension; $p_v$ is the vapor
pressure, taken constant; $\pgz$ is the partial pressure of non-condensable gas when
$R = R_0$; $\kappa$ is the polytropic exponent of that gas ($\kappa = 1$ isothermal,
$\kappa = 1.4$ adiabatic air); and $\pinf(t)$ is the prescribed far-field liquid pressure.
Table~\ref{tab:params} lists values.

\paragraph{Units check.} Multiplying Eq.~\eqref{eq:rp} by $\rho$, every term should be a
pressure. On the left, $\rho R\ddot R$ and $\rho\dot R^2$ have units
$(\mathrm{kg/m^3})(\mathrm{m^2/s^2}) = \mathrm{Pa}$. On the right, $p_v$, $\pgz$ and
$\pinf$ are pressures (the factor $(R_0/R)^{3\kappa}$ is dimensionless),
$2\sigma/R$ is $(\mathrm{N/m})/\mathrm{m} = \mathrm{Pa}$, and $4\mu\dot R/R$ is
$(\mathrm{Pa\,s})(\mathrm{m/s})/\mathrm{m} = \mathrm{Pa}$. All terms are consistent.

\paragraph{Assumptions.} The bubble stays spherical; the liquid is incompressible and
infinite in extent; the gas compresses polytropically with fixed $\kappa$, with no gas
diffusion across the wall; the vapor stays at $p_v$ (evaporation and condensation are
instantaneous); and the gas pressure is uniform inside the bubble. Incompressibility is the
assumption most likely to fail: it requires the wall Mach number $|\dot R|/c$ to stay small,
with $c \approx 1482\ \mathrm{m/s}$ in water. We will report $|\dot R|_{\max}/c$ for every run
and treat cases where it approaches about $0.1$ as outside the model's range
\todo{confirm this threshold against \cite{brennen1995} or \cite{lauterborn2010}}.

\begin{table}[t]
\centering
\caption{Parameters. Liquid properties are for water at 20\,$^\circ$C
\cite{nistwebbook}. ``Fixed'' parameters are held at the baseline except when a term is
switched off for Q3.}
\label{tab:params}
\begin{tabular}{@{}lllll@{}}
\toprule
Symbol & Meaning & Units & Baseline & Range swept \\
\midrule
$\rho$ & liquid density & $\mathrm{kg/m^3}$ & $998.2$ & fixed \\
$\mu$ & dynamic viscosity & $\mathrm{Pa\,s}$ & $1.002\times10^{-3}$ & fixed, or $0$ \\
$\sigma$ & surface tension & $\mathrm{N/m}$ & $0.0728$ & fixed, or $0$ \\
$p_v$ & vapor pressure & $\mathrm{Pa}$ & $2.34\times10^{3}$ & fixed \\
$p_0$ & ambient pressure before the event & $\mathrm{Pa}$ & $1.01325\times10^{5}$ & fixed \\
$\kappa$ & polytropic exponent & -- & $1.4$ & $1.0$ and $1.4$ \\
$R_0$ & initial radius & $\mathrm{m}$ & $100\ \mu\mathrm{m}$ (B), $10\ \mu\mathrm{m}$ (C) & $1\ \mu\mathrm{m}$ to $1\ \mathrm{mm}$ \\
$\Pi = p_1/p_0$ & pressure ratio, Case B & -- & $5$ & $1.5$ to $20$ \\
$\varepsilon$ & forcing amplitude, Case C & -- & $0.01$ & $0.01$ to $0.8$ \\
$\omega/\omega_0$ & forcing frequency, Case C & -- & -- & $0.2$ to $3$ \\
\bottomrule
\end{tabular}
\end{table}

\subsection{Test cases}\label{sec:cases}

A bubble that starts in equilibrium under constant ambient pressure never moves, so each
case below is built to produce real dynamics: an empty cavity, a pressure step, or
oscillating forcing.

\paragraph{Case A: empty-cavity (Rayleigh) collapse.} Set $\pgz = 0$, $p_v = 0$,
$\sigma = \mu = 0$ and $\pinf = \Delta p$ constant. The bubble starts at rest with nothing
inside it, so it collapses. Multiplying Eq.~\eqref{eq:rp} by $2R^2\dot R$ and integrating
once gives
\begin{equation}\label{eq:rayleigh-energy}
  \dot R^2 = \frac{2\Delta p}{3\rho}\left(\frac{R_0^3}{R^3} - 1\right),
\end{equation}
and integrating again gives the exact time to reach any radius $R$,
\begin{equation}\label{eq:rayleigh-t}
  t(R) = R_0\sqrt{\frac{3\rho}{2\Delta p}}\int_{R/R_0}^{1}\frac{x^{3/2}}{\sqrt{1 - x^3}}\,dx .
\end{equation}
Setting $R = 0$ and substituting $u = x^3$ turns the integral into a Beta function, so the
collapse time is
\begin{equation}\label{eq:tc}
  t_c = \frac{1}{3}\sqrt{\frac{3}{2}}\,B\!\left(\tfrac{5}{6},\tfrac{1}{2}\right) R_0\sqrt{\frac{\rho}{\Delta p}}
      \approx 0.91468\,R_0\sqrt{\frac{\rho}{\Delta p}},
\end{equation}
which matches the constant $0.915$ given by Rayleigh \cite{rayleigh1917} and Brennen
\cite{brennen1995}. For $R_0 = 1\ \mathrm{mm}$ and $\Delta p = 101.325\ \mathrm{kPa}$,
$R_0\sqrt{\rho/\Delta p} = 99.3\ \mu\mathrm{s}$ and $t_c \approx 90.8\ \mu\mathrm{s}$.
Equation~\eqref{eq:rayleigh-t} gives a reference value of $R$ at any time, not just $t_c$, so
it checks the whole trajectory. This case is our main correctness check. It is also a hard
test for a time integrator, because the solution is singular at $t_c$
($\dot R \to -\infty$ as $R \to 0$), so we compare up to a small final radius rather than all
the way to $t_c$.

\paragraph{Case B: gas bubble after a pressure step.} The bubble sits in equilibrium at
$\pinf = p_0$, so $\pgz$ follows from Eq.~\eqref{eq:pg0} with $p_0$. At $t = 0$ the
far-field pressure jumps to $p_1 = \Pi p_0$ and stays there:
\begin{equation}
  \pinf(t) = p_1 \quad \text{for } t > 0 .
\end{equation}
This mimics a nucleus convected out of the low-pressure zone near a blade into the recovery
region. At $t = 0^+$ the wall acceleration is $\ddot R = (p_0 - p_1)/(\rho R_0) < 0$, so the
bubble starts collapsing at once. The gas inside stiffens rapidly as $R$ drops, stops the
collapse at a finite $R_{\min}$, and the bubble rebounds and oscillates. The right-hand side
is smooth for $t > 0$, which makes this a good case for convergence studies. With
$\mu = 0$, there is an exact conservation law to check against: the energy
\begin{equation}\label{eq:energy}
  E = 2\pi\rho R^3\dot R^2 + \frac{4\pi}{3}(p_1 - p_v)R^3 + 4\pi\sigma R^2
      + \frac{4\pi}{3}\,\frac{\pgz R_0^{3\kappa} R^{3(1-\kappa)}}{\kappa - 1}
\end{equation}
(liquid kinetic energy, work against the far field, surface energy and gas internal energy)
is constant in time for $\kappa \ne 1$. Differentiating $E$ and using Eq.~\eqref{eq:rp}
confirms $dE/dt = 0$. Case B is also the workhorse for Q1 and Q3.

\paragraph{Case C: oscillating ambient pressure.} The bubble starts in equilibrium at $p_0$,
and the far-field pressure oscillates:
\begin{equation}
  \pinf(t) = p_0\left[1 - \varepsilon\sin(\omega t)\right].
\end{equation}
Linearizing Eq.~\eqref{eq:rp} about $R_0$ with $R = R_0(1 + x)$, $|x| \ll 1$, gives a
damped, driven oscillator
\begin{equation}\label{eq:linear}
  \ddot x + 2\beta\dot x + \omega_0^2 x = \frac{\varepsilon p_0}{\rho R_0^2}\sin(\omega t),
  \qquad
  \omega_0^2 = \frac{1}{\rho R_0^2}\left(3\kappa\pgz - \frac{2\sigma}{R_0}\right),
  \qquad
  \beta = \frac{2\mu}{\rho R_0^2}.
\end{equation}
With $\sigma = p_v = 0$ this reduces to the Minnaert frequency
$\omega_0 = (1/R_0)\sqrt{3\kappa p_0/\rho}$ \cite{minnaert1933}, about
$f_0 R_0 \approx 3.29\ \mathrm{m/s}$ for air in water at one atmosphere. For our baseline
$R_0 = 10\ \mu\mathrm{m}$ with all terms kept, $f_0 = \omega_0/2\pi \approx 342\ \mathrm{kHz}$
for $\kappa = 1.4$ ($288\ \mathrm{kHz}$ for $\kappa = 1$), and the quality factor is
$\omega_0/(2\beta) \approx 54$. The steady linear amplitude,
$|x| = (\varepsilon p_0/\rho R_0^2)\big/\sqrt{(\omega_0^2 - \omega^2)^2 + 4\beta^2\omega^2}$,
is an exact reference at small $\varepsilon$. Increasing $\varepsilon$ takes us into the
nonlinear regime where the peak shifts and the response can become violent, which is Q2.

\subsection{Dimensionless form}\label{sec:nondim}

Scale the radius by $R_0$ and time by the Rayleigh time $t_R$,
\begin{equation}
  \tilde R = \frac{R}{R_0}, \qquad \tilde t = \frac{t}{t_R}, \qquad
  t_R = R_0\sqrt{\frac{\rho}{\Delta p}},
\end{equation}
where $\Delta p$ is the pressure difference driving the motion: $\Delta p = p_1 - p_v$ for
Case B ($\Delta p = \pinf$ for Case A, where $p_v = 0$), and $\Delta p = p_0 - p_v$ for
Case C. With primes for $d/d\tilde t$, Eq.~\eqref{eq:rp} becomes
\begin{equation}\label{eq:rp-nondim}
  \tilde R\tilde R'' + \frac{3}{2}\tilde R'^2
  = P_g\,\tilde R^{-3\kappa} - P(\tilde t) - \frac{S}{\tilde R}
    - \frac{4}{\mathrm{Re}}\,\frac{\tilde R'}{\tilde R},
  \qquad \tilde R(0) = 1,\quad \tilde R'(0) = 0,
\end{equation}
with the dimensionless groups
\begin{equation}
  P_g = \frac{\pgz}{\Delta p}, \qquad
  P(\tilde t) = \frac{\pinf(t) - p_v}{\Delta p}, \qquad
  S = \frac{2\sigma}{R_0\,\Delta p}, \qquad
  \mathrm{Re} = \frac{R_0\sqrt{\rho\,\Delta p}}{\mu},
\end{equation}
plus $\kappa$, and for Case C the forcing amplitude $\varepsilon$ and frequency
$\omega t_R$ (or equivalently $\omega/\omega_0$, since $(\omega_0 t_R)^2 = 3\kappa P_g - S$).
In Case A, $P_g = S = 1/\mathrm{Re} = 0$ and $P = 1$, so the dimensionless collapse time is
the pure number $0.91468$. In Case B, $P = 1$ and
$P_g = (p_0 - p_v)/(p_1 - p_v) + S \approx 1/\Pi$, so the pressure ratio, $\kappa$, $S$ and
$\mathrm{Re}$ fully determine $\tilde R(\tilde t)$.

For water at one atmosphere, $S = 1$ at $R_0 \approx 1.5\ \mu\mathrm{m}$ and
$\mathrm{Re} = 1$ at $R_0 \approx 0.1\ \mu\mathrm{m}$. For $R_0 \gtrsim 10\ \mu\mathrm{m}$ both
corrections are small at the scale of the initial bubble ($S \approx 0.15$,
$\mathrm{Re} \approx 100$ at $10\ \mu\mathrm{m}$), so Case B results for different bubble sizes
should fall onto one curve when plotted against $\Pi$ in these variables, and departures
from that curve at small $R_0$ show where surface tension and viscosity take over. One
caution for Q3: these groups use the initial radius and the bulk pressure. Near $R_{\min}$
the local length scale is much smaller and the local speed much larger, so viscosity and
surface tension can matter more at the instant of collapse than $S$ and $\mathrm{Re}$ suggest.

\subsection{Why we need numerical methods}\label{sec:whynumerics}

Equation~\eqref{eq:rp} is a nonlinear second-order ODE with no closed-form solution once gas
pressure, surface tension or viscosity are included, or once the ambient pressure varies in
time. The only exact solution is the Case A limit, and even there $R(t)$ is known only
implicitly through Eq.~\eqref{eq:rayleigh-t}. The linearization in Eq.~\eqref{eq:linear}
holds only for tiny oscillations, which is exactly the regime our questions are not about:
Q1 asks for $R_{\min}$ and peak pressures in collapses where $R$ changes by a large factor,
and Q2 asks where the linear prediction stops working. So we need to integrate the IVP
numerically, across hundreds of parameter combinations.

The problem is also hard for a time integrator in a specific way. Near the end of a
collapse, Eq.~\eqref{eq:rayleigh-energy} gives $|\dot R| \propto R^{-3/2}$, so
$R \propto (t_c - t)^{2/5}$, and the local time scale $R/|\dot R| = \tfrac{5}{2}(t_c - t)$
shrinks to zero. In a gas-filled bubble the collapse is stopped at $R_{\min}$, but the time
spent near the minimum is still a small fraction of the overall period, and it gets smaller
as $\Pi$ grows. A fixed step small enough to resolve the minimum wastes most of its work
on the slow parts of the motion, which is why we compare fixed-step RK4 with an adaptive
embedded Runge--Kutta pair that can take small steps only where they are needed.
